\documentclass[12pt,letterpaper]{article}

%Packages
\usepackage{amsmath}
\usepackage{amsthm}
\usepackage{amsfonts}
\usepackage{amssymb}
\usepackage{amscd}
\usepackage{dsfont}
\usepackage{enumerate}
\usepackage{fancyhdr}
\usepackage{mathrsfs}
\usepackage{bbm}
\usepackage{framed}
\usepackage{mdframed}
\usepackage{cancel}
\usepackage{float}
\usepackage{mathtools}
%\usepackage[]{mcode}
\usepackage{graphicx}
\usepackage{tikz}
\usepackage{hyperref}
\hypersetup{
    colorlinks=true,
    linkcolor=blue,
    filecolor=magenta,      
    urlcolor=cyan,
}
\usetikzlibrary{automata,arrows}

%Page formatting
\usepackage[letterpaper,voffset=-.5in,bmargin=4cm,footskip=1cm]{geometry}
\setlength{\parindent}{0.0in}
\setlength{\parskip}{0.1in}
\allowdisplaybreaks
\headheight 15pt
\headsep 10pt

% Common Commands
\newcommand\N{\mathbb N}
\newcommand\Z{\mathbb Z}
\newcommand\R{\mathbb R}
\newcommand\Q{\mathbb Q}
\newcommand\lcm{\operatorname{lcm}}
\newcommand\setbuilder[2]{\ensuremath{\left\{#1\;\middle|\;#2\right\}}}
\newcommand\E{\operatorname{E}}
\newcommand\V{\operatorname{V}}
\newcommand\Pow{\ensuremath{\operatorname{\mathcal{P}}}}

\DeclarePairedDelimiter\ceil{\lceil}{\rceil}
\DeclarePairedDelimiter\floor{\lfloor}{\rfloor}

\newcommand\definition{\textbf{Definition: }}
\newcommand\proposition{\textbf{Proposition: }}

% 22 style
\newcommand\hint[1]{\textbf{Hint}: #1}
\newcommand\note[1]{\textbf{Note}: #1}
\newenvironment{22enumerate}{\begin{enumerate}[a.]\itemsep0em}{\end{enumerate}}
\newenvironment{22itemize}{\begin{itemize}\itemsep0em}{\end{itemize}}
\fancypagestyle{firstpagestyle} {
  \renewcommand{\headrulewidth}{0pt}%
  \lhead{\textbf{CSCI 0220}}%
  \chead{\textbf{Discrete Structures and Probability}}%
  \rhead{Littman}%
}



\pagestyle{fancyplain}

\newcommand\EX{\mathds{E}}
\newcommand\PR{\mathds{P}}
\newcommand\zlam{Z_\lambda}
\DeclareMathOperator*{\argmin}{arg\,min}

%%%%%%%%%%%%%%%% COMPILE WITH \soltrue or \solfalse %%%%%%%%%%%%%%%%%%%%%%%%%%%%%
\newif\ifsol
\soltrue  % SOLUTIONS
% \solfalse % NO SOLUTIONS
%%%%%%%%%%%%%%%%%%%%%%%%%%%%%%%%%%%%%%%%%%%%%%%%%%%%%%%%%%%%%%

%%%%%%%%%%%%%%%% solution help %%%%%%%%%%%%%%%%%%%%%
\newcommand{\solu}[2]{ \begin{mdframed} \ifsol #2 \else \vspace{#1} \fi \end{mdframed} }
\newcommand{\solt}{ \ifsol (T) \fi }
\newcommand{\solf}{ \ifsol (F) \fi }
\newcommand{\solm}[1]{\ifsol \textit{(#1)} \fi}


\begin{document}
  \thispagestyle{firstpagestyle}
  \begin{center}
    {\large \textbf{Recitation 3}}
    
    {\large Set Equivalence}

  \end{center}

\section*{Set Theory: Review}

Recall the definitions and terminology for set theory:

\textbf{Defn 1:} A \textbf{set} is a collection of objects with no repetition or order.

\textbf{Defn 2:} $B$ is a \textbf{subset} of $A$ if every element in $B$ is also in $A$. This is written as $B \subseteq A$.

\textbf{Defn 3:} $\Pow{(A)}$, called the \textbf{power set} of $A$, is the set of all subsets of $A$.

\textbf{Defn 4:} $A \cup B$ denotes the \textbf{union} of sets $A$ and $B$. This contains all the elements from $A$, and all of the elements from $B$. 
	
\textbf{Defn 5:} $A \cap B$ denotes the \textbf{intersection} of sets $A$ and $B$. This contains only the elements that appear in both of the sets.
	
\textbf{Defn 6:} $A - B$ denotes the \textbf{difference} of sets $A$ and $B$. This contains elements that appear in $A$ but not $B$. In other words, $A - B = \{x \in A \ | \ x \notin B\}$.
	
\textbf{Defn 7:} $\overline{A}$ denotes the \textbf{complement} of $A$ relative to some universal set $U$. $\bar{A} = U - A$, that is, it is everything except what is in $A$.
	
\textbf{Defn 8:} $|A|$ denotes the \textbf{cardinality} of $A$, which is a count of the number of elements contained in $A$.

\textbf{Task 1}

Define the following sets
\begin{align*}
	A &= \left\{\varnothing, 0, \left\{\varnothing\right\},
\left\{0,\varnothing\right\} \right\} \\
	B &= \left\{\varnothing, \left\{\varnothing\right\},
\left\{0, \varnothing, 1\right\} \right\} \\
	C &= \{x \mid \exists \ y \in \Z \text{ s.t. } y^2=x, \ x < 10 \} \\
    D &= \{a, b, c, d, e\} \\
    E &= \{c, d, e\} \\
    F &= \{d, e, f, g\} \\
\end{align*}
\begin{enumerate}[a.]
    \item Find the following sets:
  		\begin{enumerate}[i.]
  			\item $A \cup B$
      		\item $A \cap B$
      		\item $A - B$
      		\item $\Pow(B)$ ($=2^B$), the power set of $B$
      		\item $C - (A \cup B)$
      		\item $\{x \mid x \in B, |x| \notin C\}$
            \item $A \times E$
            \item \textit{Optional: }$(D \times E) - (E \times F)$
            \item \textit{Optional: }$(D - E) \cap (D \cap E)$
	   \end{enumerate}
  \item Find the cardinalities of the following sets:
        \begin{enumerate}[i.]
          \item $C$
          \item $A \times B$
          \item $\Pow(A \cup B)$
          \item $F - (D \cap E)$
          \item \textit{Optional: }$\Pow(\Pow(\varnothing))$
        \end{enumerate}
    \solu{5cm}{
    \begin{enumerate}
        \item \begin{enumerate}[i.]
      		\item
            $\left\{\varnothing,0,\left\{\varnothing\right\},
            \left\{0,\varnothing\right\},
            \left\{0,\varnothing,1\right\}\right\}$
      		\item
            $\left\{\varnothing,\left\{\varnothing\right\}\right\}$
      		\item
            $\left\{0, \left\{0,\varnothing\right\}\right\}$
      		\item
  				$\Big\{\hspace*{2pt} \varnothing, \left\{\varnothing\right\},
  				\left\{\left\{\varnothing\right\}\right\},
              \left\{\left\{0,\varnothing,1\right\}\right\},$

  				$\left\{\varnothing,\left\{\varnothing\right\}\right\},
              \left\{\varnothing,\left\{0,\varnothing,1\right\}\right\},$

  				$\left\{\left\{\varnothing\right\},
              \left\{0,\varnothing,1\right\}\right\},$

  				$\hspace*{7pt} \left\{\varnothing,\left\{\varnothing\right\},
  				\left\{0,\varnothing,1\right\}\right\} \hspace*{2pt}\Big\}$
      		\item
            $\{1, 4, 9\}$
      		\item
            $\left \{ \left\{ 0,\varnothing, 1 \right\} \right \}$
          \item
            $\{(\varnothing, c), (\varnothing, d), (\varnothing, e),$

               $(0, c), (0, d), (0, e),$

               $(\{\varnothing\}, c),
               (\{\varnothing\}, d), (\{\varnothing\}, e),$

               $(\{0, \varnothing\}, c),
               (\{0, \varnothing\}, d), (\{0, \varnothing\}, e)\}$
          \item
            $\{(a,c),(a,d),(a,e),(b,c),(b,d),(b,e),(c,c),(d,c),(e,c)\}$
          \item
            $\varnothing$
      	\end{enumerate}
        \item \begin{enumerate}[i.]
          \item 4
          \item 12
          \item 32
          \item 2
          \item 2
        \end{enumerate}
    \end{enumerate}}
\end{enumerate}

\textbf{Checkpoint 1 — Call over a TA!}

\newpage
\section*{Set Equivalence}
\subsection*{Set Element Method}
\textbf{Defn 1:} Two sets $S$ and $T$ are equal if and only if they have the same elements:
\begin{equation*}
    S = T \iff (\forall x : x \in S \iff x \in T)
\end{equation*}
\textbf{Defn 2:} S and T are equal if and only if both S is a subset of T and T is a subset of S.

From these two definitions, we can construct our set-element method for proving set equalities; that is, we show that two sets are equal if and only if every element of $S$ is also an element of $T$ and every element of $T$ is also an element of $S$.

In other words, to prove that one set is a subset of the other, we can show that $S \subseteq T$ for arbitrary sets $S$ and $T$ using the following steps:
\begin{enumerate}[1.]
    \item Let x be an element of S.
    \item Prove that x is an element of T.
    \item Conclude that $S \subseteq T$.
\end{enumerate}

\textbf{Example}

\textit{Claim: } $A \cap (A \cup B) = A$.

\textit{Proof: } We show that both $A \cap (A \cup B) \subseteq A$ and $A \subseteq A \cap (A \cup B)$.
\begin{enumerate}
    \item We will first prove our sub-claim that $A \cap (A \cup B) \subseteq A$.
    
    Consider any $x \in A \cap (A \cup B)$. This means that $x \in A$ and $x \in A \cup B$. In particular, we know that $x \in A$. Thus, we have shown that $A \cap (A \cup B) \subseteq A$.
    
    \item We will then prove the second sub-claim that $A \subseteq A \cap (A \cup B)$.
    
    Consider any $x \in A$. Then, $x \in A \cup B$. Thus, we know that $x \in A$ and $x \in A \cup B$, which can be rewritten as $x \in A \cap (A \cup B$. This shows that $A \subseteq A \cap (A \cup B)$.
\end{enumerate}
Therefore, $A = A \cap (A \cup B)$.

Now, let's practice!

\newpage
\textbf{Task 2}

For each of these statements, either prove (using the 'set element' method) or disprove the claim.

Hint: Use a counterexample to disprove a claim!
\begin{enumerate}[a.]
    \item For any two sets $A$ and $B$, $\Pow(A \cap B) = \Pow(A) \cap \Pow(B)$. 
    
    \solu{4cm}{\textit{Claim}: For all sets $A$ and $B$, 
        $\Pow(A \cap B) = \Pow(A) \cap \Pow(B)$.

        \textit{Proof}: 

        We will use the set element method to prove our claim. 
        We will first show $\Pow(A \cap B) \subseteq \Pow(A) 
        \cap \Pow(B)$, and we will then show $\Pow(A) \cap 
        \Pow(B) \subseteq \Pow(A \cap B)$. 

        Consider some $x$ in $\Pow(A \cap B)$. If $x$ is in 
        $\Pow(A \cap B)$, then $x$ must be a subset of $A \cap B$. By the definition of
        intersection, $x$ is a subset of $A$ and $x$ is a subset of $B$. This means $x$ must be a member of
        $P(A)$ and $x$ must be a member of $P(B)$. Therefore, $x$ must be a member of $P(A) \cap P(B)$. Thus, 
        $\Pow(A \cap B) \subseteq \Pow(A) \cap \Pow(B)$. 

        Now, consider some $y$ in $\Pow(A) \cap P(B)$. 
        By the definition of intersection, $y$ is in
         both $\Pow(A)$ and $\Pow(B)$. If $y$ is in both
          $\Pow(A)$ and $\Pow(B)$, then $y$ is a subset of 
          both $A$ and $B$. Thus, $\Pow(A \cap B)$ must also contain $y$. 
          Thus, $\Pow(A) \cap P(B) \subseteq \Pow(A \cap B)$. 

        Because $\Pow(A \cap B) \subseteq \Pow(A) \cap \Pow(B)$ and 
        $\Pow(A) \cap P(B) \subseteq \Pow(A \cap B)$, 
        $\Pow(A \cap B) = \Pow(A) \cap P(B)$.}
        

    \item For any two sets $A$ and $B$, $\Pow(A \cup B) = \Pow(A) \cup \Pow(B)$.
    
    \solu{4cm}{
    
    To disprove our claim, it suffices to provide a counterexample. 
        Let $A$ be $\{1\}$. Let $B$ be $\{3\}$. 


        $\Pow(A \cup B) = \{ \{1\}, \{3\}, \{1, 3\}, \varnothing\}$. 


        $\Pow(A) = \{ \{1\}, \varnothing\}$. 

        $\Pow(B) = \{ \{3\}, 
        \varnothing\}$.

        Therefore, $\Pow(A) \cup \Pow(B) = \{ \{1\},
         \{3\}, \varnothing\}$, which is not $\{ \{1\}, \{3\}, \{1, 3\},
          \varnothing \}$. Thus $\Pow(A \cup B)$ is not equal to $\Pow(A)
           \cup \Pow(B)$.}
    \end{enumerate}

\textbf{Task 3}

\begin{enumerate}[a.]
    \item Let $A$ and $B$ be subsets of some universal set. Then $A - (A - B) = A \cap B$.
    \solu{4cm}{Let $A$ and $B$ be subsets of some universal set. We will prove that $A - (A - B) = A \cap B$ by proving that $A - (A - B) \subseteq A \cap B$ and that $A \cap B \subseteq A - (A - B)$.
    
    \begin{enumerate}[1.]
        \item We will show that $A - (A - B) \subseteq A \cap B$.
        
        Let $x \in A - (A - B)$. This means that $x \in A$ and $x \notin (A - B)$.
    
        We know that an element is in $(A - B)$ if an only if it is in $A$ and not in $B$. Since $x \notin (A - B)$, we conclude that $x \notin A$ or $x \in B$. However, we also know that $x \in A$ and so we conclude that $x \in B$. This proves that $x \in A$ and $x \in B$.
    
        This means that $x \in A \cap B$, and hence we have proved that $A - (A - B) \subseteq A \cap B$.
    
        \item We will show that $A \cap B \subseteq A - (A - B)$.
        
        Now, we choose $y \in A \cap B$. This means that $y \in A$ and $y \in B$.
    
         We note that $y \in (A - B)$ if and only if $y \in A$ and $y \notin B$ and hence, y $\in (A - B)$ if and only if $y \notin A or y \in B$. Since we have proved that $y \in B$, we conclude that $y \notin (A - B)$, and hence, we have established that $y \in A$ and $y \notin (A - B)$. This proves that if $y \in A \cap B$, then $y \in A - (A - B)$ and hence, $A \cap B \subseteq A - (A - B)$.
    \end{enumerate}
    Since we have proved that $A - (A - B) \subseteq A \cap B$ and $A \cap B \subseteq A - (A - B)$, we conclude that $A - (A - B) = A \cap B$.}
    
    \item \textit{Optional: }Let $A$ and $B$ be the following sets:
    \begin{equation*}
    \begin{aligned}
        A =& \{6n : n \in \mathbb{Z}\} \\
        B =& \{2n : n \in \mathbb{Z}\} \cap \{3n :  n \in \mathbb{Z}\}
    \end{aligned}
    \end{equation*}
    Prove that $A = B$.
    
    \solu{4cm}{
    We will show that both $A \subseteq B$ and $B \subseteq A$.
    
    \begin{enumerate}[1.]
        \item First, we prove that $A \subseteq B$; that is, $\{6n$: n $\in \mathbb{Z}\} \subseteq \{2n: n \in \mathbb{Z}\} \cap \{3n$ : $ n \in \mathbb{Z}\}$.
        
        Choose any $m \in \{6n$: $n \in \mathbb{Z}\}$. By definition, we can write $m = 6k$ for some $k \in \mathbb{Z}$.
        
        In particular, $m = 2(3k)$, so that $m \in \{3n$: $n \in \mathbb{Z}\}$ as well. By the definition of intersection, we can conclude that $m \in \{2n$: $n \in \mathbb{Z}\} \cap \{3n$: $n \in \mathbb{Z}\}$ as required.
        
        \item Next, we prove that $\{2n: n \in \mathbb{Z}\} \cap \{3n$ : $ n \in \mathbb{Z}\} \subseteq \{6n$: n $\in \mathbb{Z}\}$.
        
        Choose any $m \in \{2n$: $n \in \mathbb{Z}\} \cap \{3n$: $n \in \mathbb{Z}\}$, so that both $m \in \{2n$: $n \in \mathbb{Z}\}$ and $m \in \{3n$: $n \in \mathbb{Z}\}$.
        
        In particular, we can find integers $j$ and $k$ such that $m = 2j$ and $m = 3k$. The fact that $m = 2j$ means that $m$ is even, and so $m = 3k$ is even. But the product of two integers is odd if and only if both intgers are odd; since $4$ is odd but $3k$ is even, we conclude that k is even. Then, $k= 2a$ for some integer $a$, and so $m = 3k = 3(2a) = 6a$, which shows that $m \in \{6n$:$n \in \mathbb{Z}\}$ as required.
    \end{enumerate}
    
    Therefore, we can conclude that $A=B$ as we have shown both $A \subseteq B$ and $B \subseteq A$.}
    
    \item \textit{Optional:} Prove that $A \times (B \cup C) = (A \times B) \cup (A \times C)$.
    
    \solu{4cm}{
    \begin{enumerate}[1.]
        \item We will show that $A \times (B \cup C) \subseteq (A \times B) \cup (A \times C)$.
        
        Let $(x, y) \in A \times (B \cup C)$. Then $x \in A$ and $y \in B \cup C$. So $x \in A$ and $y \in B$, or $x \in A$ and $y \in C$. In the former case, $(x, y) \in A \times B$, and in the latter case $(x, y) \in A \times C$. So $(x, y) \in (A \times B) \cup (A \times C)$. Since $(x, y)$ was chosen arbitrarily, $A \times (B \cup C) \subseteq (A \times B) \cup (A \times C)$.
        
        \item We will show that $(A \times B) \cup (A \times C) \subseteq A \times (B \cup C)$.
        
        Let $(x, y) \in (A \times B) \cup (A \times C)$. Then $(x, y) \in A \times B$ or $(x, y) \in A \times C$. So $x \in A$ and $y \in B \cup C$, and hence $(x, y) \in A \times (B \cup C)$. Since $(x, y)$ was chosen arbitrarily, $(A \times B) \cup (A \times C) \subseteq A \times (B \cup C)$.
    Therefore, we can conclude that $A \times (B \cup C) = (A \times B) \cup (A \times C)$.
    \end{enumerate}}
\end{enumerate}

\textbf{Checkpoint 2 — Call over a TA!}

\newpage
\subsection*{Set Algebra}
We can also prove that two sets are equivalent using set algebra. More concretely, we can use the stated laws of set algebra to convert one side of the equation to the other (or convert both sides to an identical expression).

Example (from the sample proofs on our website):

Prove that $(A \cap B) \cup (A - B) = A \cap (B \cup (A - B))$.
\begin{flalign*}
& (A \cap B) \cup (A - B) \\
&= (A \cap B) \cup (A \cap \overline{B}) && \text{(Set Difference Law)} \\
&= (A \cap (B \cup \overline{B}) && \text{(Distribution)} \\
&= A \cap U && \text{(Complement Law)} \\
&= A && \text{(Identity Law)} \\
&= A \cap (A \cup B) && \text{(Absorption)} \\
&= A \cap (B \cup A) && \text{(Commutivity)} \\
&= A \cap ((B \cup A) \cap U) && \text{(Identity Law)} \\
&= A \cap ((B \cup A) \cap (B \cup \overline{B}) && \text{(Complement Law)} \\
&= A \cap (B \cup (A \cap \overline{B})) && \text{(Distribution)} \\
&= A \cap (B \cup (A - B)) && \text{(Set Difference Law)}
\end{flalign*}
Therefore, the equality holds since all these steps are biconditionally true.

\textbf{Task 4}
\begin{enumerate}[a.]
    \item Let $A$ and $B$ be subsets of some universal set $U$. Prove that $(A \cap \overline{B}) \cup B = A \cup B$ using set algebra.
    \solu{4cm}{\begin{flalign*}
& (A \cap \overline{B}) \cup B \\
&= (A \cup B) \cap (\overline{B} \cup B) && \text{(Distributive Law)}\\
&= (A \cup B) \cap U && \text{(Complement Law)}\\
&= A \cup B && \text{(Identity Law)}
\end{flalign*}}

    \item Let $A$ and $B$ be subsets of some universal set $U$. Prove that $(A - B) \cup (B - A) = (A \cup B) - (A \cap B)$ using set algebra.
    
    \solu{4cm}{\begin{flalign*}
& (A - B) \cup (B - A) \\
&= (A \cap \overline{B}) \cup (B \cap \overline{A}) && \text{(Set Difference Law)}\\
&= (A \cap \overline{B} \cup B) \cap ((A \cap \overline{B} \cup \overline{A}) && \text{(Distributive Law)}\\
&= ((A \cup B) \cap (\overline{B} \cup B)) \cap ((A \cup \overline{A}) \cap (\overline{B} \cup \overline{A})) && \text{(Distributive Law)}\\
&= ((A \cup B) \cap U) \cap (U \cap \overline{B} \cup \overline{A}) && \text{(Complement Law)}\\
&= (A \cup B) \cap (\overline{B} \cup \overline{A}) && \text{(Identity Law)}\\
&= (A \cup B) \cap (\overline{A \cap B}) && \text{(DeMorgan's Law)}\\
&= (A \cup B) - (A \cap B) && \text{(Set Difference Law)}
\end{flalign*}}

    \item \textit{Optional: }Show that $(A - B) - (B - C) = A - B$.
    
    \solu{4cm}{\begin{flalign*}
& (A - B) - (B - C) \\
&= (A \cap \overline{B}) \cap (\overline{B \cap \overline{B}}) && \text{(Set Difference Law)}\\
&= (A \cap \overline{B}) \cap (\overline{B} \cup C) && \text{(DeMorgan's Law)}\\
&= ((A \cap \overline{B}) \cap \overline{B}) \cup ((A \cap \overline{B}) \cap C) && \text{(Distributive Law)}\\
&= (A \cap (\overline{B} \cap \overline{B})) \cup (A \cap (\overline{B} \cap C)) && \text{(Associative Law)}\\
&= (A \cap \overline{B}) \cup (A \cap (\overline{B} \cap C)) && \text{(Idempotent Law)}\\
&= A \cap (\overline{B}) \cup (\overline{B} \cap C)) && \text{(Distributive Law)}\\
&= A \cap \overline{B} && \text{(Absorption Law)}\\
&= A - B && \text{(Set Difference Law)}\\
&= (A \cap \overline{B} \cup B) \cap ((A \cap \overline{B} \cup \overline{A}) && \text{(Distributive Law)}\\
&= ((A \cup B) \cap (\overline{B} \cup B)) \cap ((A \cup \overline{A}) \cap (\overline{B} \cup \overline{A})) && \text{(Distributive Law)}\\
&= ((A \cup B) \cap U) \cap (U \cap \overline{B} \cup \overline{A}) && \text{(Complement Law)}\\
&= (A \cup B) \cap (\overline{B} \cup \overline{A}) && \text{(Identity Law)}\\
&= (A \cup B) \cap (\overline{A \cap B}) && \text{(DeMorgan's Law)}\\
&= (A \cup B) - (A \cap B) && \text{(Set Difference Law)}
\end{flalign*}}

\end{enumerate}

\newpage
\section*{Google Ad Personalization}
\textbf{Task 5}
 \begin{enumerate}[a.]
    \item Read \href{https://www.practicalecommerce.com/how-google-collects-data-to-personalize-ads}{this article} on Google's ad personalization, and how it collects data on its users to select targeted ads. 
    
    \item If you have a Google account that you actively use, follow the instructions on \href{https://www.businessinsider.com/what-does-google-know-about-me-search-history-delete-2019-10#first-navigate-to-your-google-account-homepage-by-clicking-the-widget-in-the-top-right-corner-of-any-google-site-1}{this page} and look at the groups that Google believes you are a part of. If you have ad penalization enabled and feel comfortable doing so, skim your entire ad profile. Did anything surprise you? Why or why not?
    
    \solu{}{
     \textit{Pass condition:} Students should be able to identify (if they opted to look at their profile, which is optional) certain groups or trends Google believes they are a part of. Any sort of reaction is fine, as it is impossible to predict how many groups google has put them.
     }
    
    \item 
    \textbf Google defines your ad profile as the intersection of several traits, or sets. This profile will determine what ads you are shown - advertisers choose what demographic groups and interest groups they want their ads to target (you can learn more about interest targeting \href{https://support.google.com/google-ads/answer/2497941}{here} and demographic targeting \href{https://support.google.com/google-ads/answer/2580383}{here}). Think about the assumptions and biases that come with putting people into sets. What problems and inaccuracies could arise by defining what content individuals receive solely based on the sets they belong to?
    
    \solu{}{
     \textit{Pass condition:} Students should be able to identify at least one disadvantage with defining people's content consumption by wide sets, such as stereotyping, inaccuracy, or confirmation bias. They should be able to understand that not every member of a set is the same, but designing programs around what a set "believes" runs the risk of treating every element in that set as interchangeable.
     
     \textit{Examples:} 
     \begin{itemize}
        \item I think being able to define which demographic ads are marketed towards could be really dangerous in terms of stereotyping. By defining your product to be a "female" or "male" product it can perpetuate stereotypes about gender roles, like advertising barbies to females and Nerf guns to males. This is not to mention the dangerous "other" category which binds all non-conforming identities together.
    	\item I think it's unfair to assume that 18-34 year old's all have the same interests, it's kinda weird that its a trait people can market towards. Grouping people by sets doesn't really leave any nuance between people, as the advertiser will treat every 18-34 year old as if they were the same person.
    	\item Using groups to determine ads seems weird to me. Google's profile just have some weird inaccuracies, and basing content solely on these inaccurate profiles seems like Google would need a lot of data to collect to get accurate results. This scares me since I can't control how much data google gets from my searches and web history.
     \end{itemize}


    \textit{Prodding questions:}
    \begin{22itemize}
        \item Let's say an advertiser wants to market to middle-class people for their product. Would it be fair to say all middle-class people have similar tastes and interests?
        \item What problems could arise by branding a product, like a toy, to be solely marketable to "males" and "females"?
        \item When an advertiser states they want to market to a certain group, they necessarily assume people in that group have the same taste. How might this gloss over differences between members in that group?
    \end{22itemize}
     }
    
    \item Being at the intersection of several traits means that people's experiences could differ drastically from people in one set or the other. For example, take two traits that tend to be underprivileged: low education level and low income level. How might a upper-class high school dropout's opportunities and livelihood differ from a high school dropout experiencing poverty? What are some other cases of people who belong to the same set in one instance but experience vastly different amounts of privilege and discrimination based on what other sets they are a member of?
    
    % Is describing their experience as simple as combining the experience of one group and the experience of the other group, or is there something unique about being at the intersection of several groups?
    % What are other cases where people who both identify with one set might experience privilege and discrimination differently because of the existence of intersections?
    
    \solu{}{
    \textit{Pass condition:} Students should be able to recognize and identify the differences in the situations between the two people in the given example (upper-class high school dropout vs. high school dropout experiencing poverty). They should be able to draw at least two comparisons between people who share a common set(s) but potentially have different life experiences, and students should be able to explain why they chose these comparisons. They should understand the core idea of intersectionality - being a member of two disadvantaged groups has unique struggles. For example, discussing discrimination against black women is not as easy as branding things as either "racism" or "sexism"; their experience is unique.
    
    \textit{Examples:}
    Identifying potential differences between two high-school dropouts, when one is upper-class and the other is in poverty:
    \begin{22itemize}
        \item Upper-class high school dropout could potentially have financial support from their family, while a dropout experiencing poverty may not have the same luxury
        \item The two may have had different opportunities before they dropped out, leading to them possessing different skill sets. 
    \end{22itemize}
    Comparisons between people who may share a set but have different experiences in life:
    \begin{22itemize}
        \item White man vs. white woman in STEM (discrimination towards women in STEM, belief that women are not as capable)
        \item Black vs white woman
        \item Any situation in which the student lists a combination of two social or personal identities and points out discrepancies between being being part of only one or two of the listed groups
    \end{22itemize}
    
    \textit{Prodding questions:}
    \begin{22itemize}
        \item Why might an upper-class dropout be more keen on taking financial risks such as investing in the stock market or entrepreneurship?
        \item Would an upper-class
        \item Think about historically disadvantaged groups. Why is being a member of two of these groups uniquely worse than being  member of only one?
    \end{22itemize}
    }
    
    \item When it comes to data analysis and storage, it is common practice to group people and profiles into discrete sets such as race, gender identity, and geographical location. Standard database design requires people to be grouped together based on certain traits, with each trait being represented by a column. What potential conflicts and problems could occur when trying to fit people into a finite number of sets? How could assumptions that apply to the entire set harm the individuals in them?
    
    \solu{}{
    \textit{Pass condition:} Students should be able to recognize that sets can be extremely specific, which means that there are an infinite number of sets that people can belong to; because each person's experiences are wholly unique to them, they may belong to many more sets than the finite number of columns in a data set. Students should also be able to recognize that people are not wholly defined by one set; making assumptions and applying them broadly to everyone in the set is stereotyping.
    
    \textit{Prodding questions:} 
    \begin{22itemize}
        \item We have mostly discussed sets in terms of larger demographic groups, such as race, gender, age, etc. However, that doesn't mean sets are limited to just these broad categories (think of your Google ad profile). Are there really a finite number of sets that people can belong to?
        \item Can human experiences truly be counted?
        \item Is it possible fr people to be defined by just one characteristic? What are the pitfalls of treating others as if they are one-dimensional?
    \end{22itemize}
    }
\end{enumerate}
\subsection*{Checkoff — Call over a TA!}

\end{document}
